% Simple Beamer presentation for Project_1
% Compile with: `pdflatex presentation.tex` or `xelatex presentation.tex`
\documentclass{beamer}
\usepackage[utf8]{inputenc}
\usepackage{graphicx}
\usepackage{amsmath}
\usepackage{xcolor}
\usetheme{Madrid}

% Dark (black) theme settings
\definecolor{bgcolor}{RGB}{255,255,255}
\definecolor{fgcolor}{RGB}{0,0,0}
\definecolor{accent}{RGB}{255,165,0}
\setbeamercolor{background canvas}{bg=bgcolor}
\setbeamercolor{normal text}{fg=fgcolor,bg=bgcolor}
\setbeamercolor{frametitle}{fg=fgcolor,bg=bgcolor}
\setbeamercolor{title}{fg=fgcolor}
\setbeamercolor{author}{fg=fgcolor}
\setbeamercolor{institute}{fg=fgcolor}
\setbeamercolor{section in toc}{fg=fgcolor}
\setbeamercolor{itemize item}{fg=accent}
\setbeamercolor{enumerate item}{fg=accent}

% Bold and uppercase frame titles to match uploaded presentation
\setbeamerfont{frametitle}{series=\bfseries}

\title{\textbf{Modeling 2D Marshak Waves}}
\author{Yuval Kochavi \\ Dr. Shai Heizler}
\institute{The Hebrew University of Jerusalem}
\date{17.6.2026}

\begin{document}

\begin{frame}
  \titlepage
\end{frame}

\begin{frame}{MOTIVATION}
  \begin{itemize}
    \item Cylindrical radiation diffusion problems arise in many physical scenarios.
    \item The 1D solution only provides centerline front position $z_F(t)$.
    \item Real experiments display complicated 2D geometry requiring rigorous modeling. 
  \end{itemize}
\end{frame}

\begin{frame}{THE MECHANISM}
  \begin{itemize}
    \item Radiation propagates axially (along $z$) and radially (along $r$).
    \item Wall absorption causes overall energy loss.
    \item Radial temperature gradients dynamically develop.
    \item The radiation-driven heat front curves in $(r,z)$ space. 
  \end{itemize}
\end{frame}

\begin{frame}{PREVIOUS WORKS}
  \begin{itemize}
    \item R. E. Marshak  1958 - First description and super-sonic solutions. 
    \item J. H. Hammer; M. D. Rosen  2006 - New 1D solutions based on a perturbation theory.
    \item O. A. Hurricane; J. H. Hammer  2006 - Matematical characterization of marshak waves curvature.
    \item A. P. Cohen; G. Malamud, S. I. Heizler  2020 - Modeling the 1D heat front Based on many experiments.
  \end{itemize}
\end{frame}

\begin{frame}{OUR WORK - OVERVIEW}
  \begin{itemize}
    \item The idea: Solve the 1D propagation and radial perturbation structure separately and patch them.
    \item We deploy separation of variables on the 2D energy equation.
    \item The fundamental Bessel mode is introduced mathematically.
    \item The radial structure of the heat front elegantly surfaces as an analytical Bessel profile.
  \end{itemize}
\end{frame}

\begin{frame}{ASSUMPTIONS}
  \begin{itemize}
    \item \textbf{Geometry:} Domain spans $0 \leq r \leq R$ (radius) and $0 \leq z \leq L$ (length).
    \item \textbf{Boundary Condition at Source:} Marshak boundary condition relates incoming and outgoing radiation safely at $z=0$. 
    \item \textbf{Wall Interaction:} Parameterized via Albedo (reflectivity factor $\alpha_{\text{wall}}$).
    \item \textbf{Product Solution:} We adopt a separated solution profile: $T(r,z,t) = T_0(z,t) \cdot \Phi(r)$.
  \end{itemize}
\end{frame}

\begin{frame}{STATEMENT OF THE PROBLEM}
  The radiation diffusion equation in foam material:
  \begin{equation}
    \frac{\partial U}{\partial t} = \nabla \cdot (D \nabla U)
  \end{equation}
  where $U = a T^4$ is the radiation energy density, and $D = \frac{c\lambda}{3}$ is the radiation diffusion coefficient.

  \vspace{10pt}
  In cylindrical coordinates $(r, z)$:
  \begin{equation}
    \frac{\partial U}{\partial t} = \frac{1}{r}\frac{\partial}{\partial r}\left(r D \frac{\partial U}{\partial r}\right) + \frac{\partial}{\partial z}\left(D \frac{\partial U}{\partial z}\right)
  \end{equation}
\end{frame}

\begin{frame}{THE 1D CENTERLINE SOLUTION}
  For the self-similar heat front at the exact centerline ($r=0$):
  \begin{equation}
    x_F^2(t) = \frac{2+\epsilon}{1-\epsilon} \cdot C(\rho) \cdot H(t)^{-\epsilon} \cdot \int_0^t H(t') \, dt'
  \end{equation}
  An implicit time-marching numerical approach acts on:
  \begin{equation}
    E(t_i) = E(t_{i-1}) + \frac{1}{2}\left[F_{\text{net}}(t_i) + F_{\text{net}}(t_{i-1})\right] \Delta t
  \end{equation}
  incorporating wall energy loss $F_{\text{net}}(t_i) = F(0, t_i) - \frac{\dot{E}_{\text{wall}}(t_i)}{A}$.
\end{frame}

\begin{frame}{WALL ENERGY LOSS \& ALBEDO}
  Energy transfer to the cylindrical wall at $r=R$:
  \begin{equation}
    \frac{dE_{\text{wall}}}{dt} = 2\pi R L \sigma_{SB} \left(T_s^4 - \frac{\dot{E}_{\text{wall}}}{4\pi R L \sigma_{SB}}\right)
  \end{equation}
  The albedo $\alpha_{\text{wall}}$ tracks the non-linear relationship between reflected and incident flux:
  \begin{equation}
    \alpha_{\text{wall}} = \frac{F_{\text{in}}}{F_{\text{out}}} = \frac{\sigma_{SB}T_s^4 + \frac{\dot{E}_{\text{wall}}}{2}}{\sigma_{SB}T_s^4 - \frac{\dot{E}_{\text{wall}}}{2}}
  \end{equation}
\end{frame}

\begin{frame}{THE 2D RADIAL EIGENVALUE PROBLEM}
  Under assumed product formulation, the radial variant behaves linearly:
  \begin{equation}
    \frac{1}{r}\frac{d}{dr}\left(r\frac{d\Phi}{dr}\right) + \kappa^2 \Phi = 0
  \end{equation}
  A Robin boundary condition dictates behavior at the wall $r=R$, securing a continuous profile without divergence. Thus $\Phi(r) = J_0(\kappa r)$. 
  \vspace{10pt}

  The boundary enforces a transcendental equation for eigenvalues $\kappa$:
  \begin{equation}
    \kappa R \cdot J_1(\kappa R) = \epsilon_{\text{loss}} \cdot J_0(\kappa R)
  \end{equation}
\end{frame}

\begin{frame}{RESULTS: CURVED FRONT POSITION}
  \textbf{Two-Dimensional Temperature Distribution}:
  The front curves precisely via the first eigenvalue $\kappa_0$:
  \begin{equation}
    z_F(r,t) = z_F(t) \cdot J_0(\kappa_0 r)
  \end{equation}
  
  \textbf{Key observations:}
  \begin{itemize}
      \item The front exhibits max penetration directly on axis ($r=0$).
      \item \textbf{Perfect Reflector ($\alpha_{\text{wall}} \to 1$):} Reduces $\kappa_0 \to 0$ generating a decoupled flat 1D plane. 
      \item \textbf{Total Absorber ($\alpha_{\text{wall}} \to 0$):} Compresses the function giving sharp decay off-axis.
  \end{itemize}
\end{frame}

\begin{frame}{RESULTS: REPRESENTATIVE SIMULATION}
  \begin{figure}
    \centering
    \IfFileExists{Figures_new/ExampleFigure.png}{%
      \includegraphics[width=0.75\textwidth]{Figures_new/ExampleFigure.png}
    }{%
      \setlength{\fboxsep}{8pt}%
      \fbox{\parbox{0.75\textwidth}{\centering Missing figure: \texttt{Figures\_new/ExampleFigure.png}\\[6pt]Place the generated image in the Figures\_new/ folder or update the path.}}
    }
    \caption{Self-Similar Solution of the 2D Heat Wave plotted against time.}
  \end{figure}
\end{frame}

\begin{frame}{COMPARISON TO NUMERICAL SIMULATIONS (1)}
  \begin{itemize}
    \item The analytically derived $J_0$ mapping maintains remarkably solid alignment with rigorous 2D simulations.
    \item Insert actual simulation plots here proving $T(r,z,t)$ validation.
  \end{itemize}
\end{frame}

\begin{frame}{COMPARISON TO NUMERICAL SIMULATIONS (2)}
  \begin{itemize}
    \item Highlight variations of albedo. 
    \item Discuss errors mapping extreme $\epsilon$ shifts.
  \end{itemize}
\end{frame}

\begin{frame}{SUMMARY}
  \begin{itemize}
    \item We offer an explicit analytical derivation coupling the 1D self-similar wave front equations onto 2D bounded spaces.
    \item Separation of variables generates a natural $J_0(\kappa_0 r)$ curvature.
    \item The model elegantly wraps radiation losses through empirical albedo adjustments across the wall.
    \item Accurate, cheap, and easily adapted for varying pulse profile shapes.
  \end{itemize}
\end{frame}

\end{document}
